\section{Computational certification and reproducibility}
\label{sec:reproducibility}

\subsection{Proof architecture}

The computational part of the proof has five layers.
\begin{enumerate}[label=(\arabic*),leftmargin=2.2em]
\item A from-scratch mask scan constructs the connected and eligible
candidate universe, canonicalises left orbits, and runs the exact-cover
decision of \cref{alg:exact-cover}.
\item Before selecting positives, the global congruence gate constructs the
rank-three corner locus on all $2{,}874$ eligible candidates, proves its
carrier covariance, and exhausts every invariant-surviving isometry in two
exact lanes.
\item Deterministic builders attach the Patterson, cover, selector, contact,
boundary, and ambient-symmetry records to the immutable $82$ row identifiers.
\item Independent algorithms cross-check the load-bearing quotient counts:
source replay versus mask scan, recursive versus combination exact cover,
orbit canonicalisation versus Burnside, and distance-graph versus affine-Gram
isometry enumeration.
\item Validators check schemas, self-hashes, row joins, theorem totals,
negative controls, and byte identity of regenerated certificates.
\end{enumerate}

No random seed, numerical tolerance, network service, or proprietary solver
is involved.  Group, mask, incidence, and flag-complex calculations use
integers; ambient coordinates, distances, and plane data use exact rational
arithmetic.

\subsection{Replay modes}

The companion repository provides a short verifier for the pinned theorem
certificate and a full reconstruction mode.  The short path verifies all
displayed totals, row identities, self-hashes, parent pins, complete joins,
and the existence of every public artifact cited in the
theorem-to-certificate crosswalk.  The full path regenerates the connected
and eligible census, the all-candidate congruence gate, and the
congruent-copy exact-cover decision, then compares all $2{,}874$ candidate
records with the canonical certificate.  The Patterson, cover, selector,
contact, boundary, and ambient-symmetry layers are frozen exact artifacts:
the public package validates their provenance, row pins, schemas, joins, and
structural identities, but does not regenerate them from their internal
production builders.  A fresh-clone gate runs the package in a temporary
directory to ensure that no neighbouring repository or private absolute
path is required.

The preflight additionally checks:
\begin{itemize}[leftmargin=2em]
\item deterministic JSON serialisation and SHA-256 manifests;
\item absence of private paths and internal routing material;
\item a completely green public test suite;
\item successful title-named LaTeX compilation with resolved citations and
cross-references;
\item agreement between every number displayed in the paper and the public
theorem certificate.
\end{itemize}

\subsection{Finite proof versus theorem import}

The exact enumeration proves statements about the finite carrier once the
objects and equivalences are fixed.  Euclidean use of the word ``convex''
additionally imports the connected three-way convexity theorem from
\cite{BabanskyyT2}.  Bibliographic background on Patterson functions,
balanced complexes, group factorisation, and tiling terminology is not used
as an oracle for any finite count.  This distinction is recorded explicitly
in the source map and claim ledger shipped with the release-candidate
repository.

\subsection{Auditability}

The complete $82$-row catalogue is printed in \cref{app:catalogue}; larger
objects, such as the $212$ contact-edge records, are supplied in canonical
machine-readable form.  \Cref{app:certificate-crosswalk} maps every theorem
to its defining data, replay command, and certificate fields.  The public
reviewer guide gives a ten-minute reading path and a longer independent replay
path.  Release status is intentionally separate from mathematical status:
the package becomes a publication only after the final external red team and
owner-controlled release gate.

\section{Conclusion}\label{sec:conclusion}

The earlier convex classification exposed a sharp nonconvex remainder of
$82$ left-translate tiling orbits.  The all-candidate rigidity theorem proves
that free congruence creates neither additional identifications nor mixed-
orbit Coxeter-pure covers, so the left census is the complete answer to the
earlier congruent-copy problem.  On that foundation, the present work turns
the nonconvex remainder into a complete finite atlas: canonical rows, an
intrinsic labelled separator, all exact covers and multiplier fibres, exact
contact complexes, and full tetrahedral symmetry data.  The labelled
Patterson vector gives the nonconvex counterpart of the convex contact
fingerprint, while the $82/20/33/32$ compression profile records why
nonabelian coordinate labels are indispensable.

The atlas also displays two independent forms of nonuniqueness.  Three tile
rows possess more than one cover orbit, and two pairs of those covers agree
edgewise without one global contact alignment.  At the same time, every
whole partition remains tile-transitive inside the tetrahedron.  These facts
show why tile, cover, selector, local contact, and global symmetry are best
treated as separate typed layers rather than compressed into one invariant.
